% =============================================================================
% Chapitre 1 — Modélisation, DDL et contraintes d'intégrité bancaire
% =============================================================================
\chapter{Modélisation, DDL et contraintes d'intégrité bancaire}

\epigraphe{Une donnée fausse qui entre dans la base n'en ressort jamais seule.}

Ce premier chapitre pose les fondations. Tu vas comprendre \textbf{comment on découpe} les données en
tables, \textbf{comment on choisit} les types et les clés, et surtout \textbf{comment on empêche} une
donnée incohérente d'entrer dans la base. Tu créeras ensuite un dossier de prêt complet, du client
jusqu'à son tableau d'amortissement, en une seule transaction.

\begin{observation}[Tes fichiers pour ce chapitre]
\begin{itemize}
\item Ton code : \code{sql/02_chapitre1.sql}. Le corrigé complet : \code{corriges/chapitre1.sql}.
\item Prérequis : les tâches 1 à 3 de la partie 5 (base créée par le schéma, puis chargée par le
script Python).
\end{itemize}
\end{observation}

\begin{situation}{Mise en situation --- Le client de 15 ans}
Le mois dernier, dans l'ancien tableur, un prêt a été accordé à un client de 15 ans : une faute de
frappe, il en avait 51. Le même mois, une mensualité a été saisie deux fois comme payée, et un dossier
refusé est apparu dans le rapport des prêts en cours. Personne ne s'en est aperçu pendant trois mois.

Ta responsable est claire : \emph{« Je ne veux pas qu'on corrige ces erreurs. Je veux qu'elles soient
impossibles. »}
\end{situation}

% =============================================================================
\section{Cours et notions importantes}

\subsection{Pourquoi plusieurs tables ? La normalisation}

\begin{imagine}
Imagine toutes les données dans un seul grand tableau, une ligne par mensualité. Pour le prêt d'Awa,
36 lignes répètent 36 fois son nom, son revenu, son adresse e-mail, le montant de son prêt et son
taux. Awa change d'adresse e-mail : il faut modifier 36 lignes. Tu en oublies une : la base contient
désormais deux adresses pour la même personne. Laquelle est la bonne ?
\end{imagine}

Ce grand tableau souffre de trois \textbf{anomalies} classiques :
\begin{itemize}
\item \textbf{de mise à jour} : une information répétée doit être modifiée partout, sinon elle devient
contradictoire ;
\item \textbf{d'insertion} : impossible d'enregistrer un client tant qu'il n'a pas de mensualité ;
\item \textbf{de suppression} : supprimer la dernière mensualité d'un prêt efface aussi le client.
\end{itemize}

\begin{definition}{Normalisation}
Normaliser, c'est découper les données en tables de sorte que \textbf{chaque fait soit stocké une seule
fois, à un seul endroit}. Les trois premières formes normales en donnent la méthode :
\begin{enumerate}
\item \textbf{1FN} : chaque cellule contient une valeur atomique (pas de liste « 25/10, 25/11, 25/12 »
dans une case) ;
\item \textbf{2FN} : chaque colonne dépend de \textbf{toute} la clé. Dans une table dont la clé est
(prêt, numéro d'échéance), le montant du prêt ne dépend que du prêt : il va dans la table des prêts ;
\item \textbf{3FN} : aucune colonne ne dépend d'une autre colonne non clé. Le revenu dépend du client,
pas du prêt : il va dans la table des clients.
\end{enumerate}
\end{definition}

C'est exactement ce découpage que tu as vu dans la partie 3 : \code{clients}, \code{demandes_pret},
\code{echeancier_remboursements}. Chaque table décrit \textbf{une seule chose}, et les clés étrangères
les relient.

\begin{observation}[Normaliser la source, dénormaliser pour l'analyse]
Pour entraîner un modèle de Machine Learning, on a besoin au contraire d'une seule grande table, une
ligne par prêt avec toutes ses variables. On la construit \emph{à partir} de la base normalisée, par
des jointures, dans une vue ou un entrepôt de données. La règle : le système qui enregistre les
opérations reste normalisé ; les tables d'analyse peuvent être dénormalisées.
\end{observation}

\subsection{Choisir les types de données}

Le type d'une colonne est la première contrainte : il refuse tout ce qui n'a pas la bonne forme.

\begin{center}\small
\begin{tabularx}{\linewidth}{@{}P{3.1cm} P{3.6cm} L@{}}
\toprule
\textbf{Type} & \textbf{Dans le projet} & \textbf{À savoir} \\
\midrule
\code{SMALLINT}, \code{INTEGER}, \code{BIGINT} & âge, durée ; identifiants & 2, 4 ou 8 octets. Les
identifiants en \code{BIGINT} : on ne veut jamais tomber à court de numéros. \\
\code{NUMERIC(p,s)} & montants, taux, PD & décimal exact : $p$ chiffres dont $s$ après la virgule.
\code{NUMERIC(12,2)} va jusqu'à 9\,999\,999\,999,99. \\
\code{REAL}, \code{DOUBLE PRECISION} & jamais pour l'argent & approximatifs : $0{,}1 + 0{,}2 \neq 0{,}3$. \\
\code{VARCHAR(n)}, \code{TEXT} & noms, statuts, descriptions & \code{VARCHAR(n)} refuse plus de $n$
caractères, \code{TEXT} n'a pas de limite. \\
\code{CHAR(1)} & note de risque & longueur fixe. \\
\code{BOOLEAN} & \code{defaut_anterieur}, \code{est_champion} & vrai, faux ou NULL. \\
\code{DATE} & échéances, paiements & un jour, sans heure. \\
\code{TIMESTAMP} & \code{date_prediction} & date et heure. En production, on préfère souvent
\code{TIMESTAMPTZ}, qui tient compte du fuseau horaire. \\
\bottomrule
\end{tabularx}
\end{center}

Les dates méritent une attention particulière, car tout le calcul du retard repose sur elles.

\begin{itemize}
\item \code{DATE - DATE} donne un \textbf{nombre entier de jours} :
\code{DATE '2026-09-30' - DATE '2026-08-15'} vaut \textbf{46}. C'est le DPD du prêt 61.
\item \code{DATE + entier} ajoute des jours ; \code{DATE + INTERVAL '1 month'} ajoute un mois et
renvoie un \code{TIMESTAMP} (d'où les conversions \code{::date}).
\end{itemize}

\begin{piege}{Piège --- Ajouter des mois en chaîne}
Que vaut le 31 janvier plus un mois ? PostgreSQL renvoie le \textbf{28 février} : il s'arrête au
dernier jour du mois. Le problème apparaît si l'on enchaîne les ajouts :
\begin{lstlisting}[style=sql]
SELECT DATE '2026-01-31' + INTERVAL '1 month'                              AS plus_1_mois,   -- 2026-02-28
       (DATE '2026-01-31' + INTERVAL '1 month') + INTERVAL '1 month'       AS chaine_2_mois, -- 2026-03-28
       DATE '2026-01-31' + 2 * INTERVAL '1 month'                          AS direct_2_mois; -- 2026-03-31
\end{lstlisting}
En partant de chaque échéance pour calculer la suivante, toutes les échéances après février seraient
décalées au 28. La bonne méthode : toujours partir de la \textbf{première} échéance et ajouter
$(k-1)$ mois. C'est ce que fait le corrigé de l'exercice 1.1.
\end{piege}

\subsection{Les clés : identifier et relier}

\begin{definition}{Clé primaire}
Une colonne (ou un groupe de colonnes) qui identifie \textbf{chaque ligne de façon unique}. Elle est
toujours \code{NOT NULL} et \code{UNIQUE}. Une table a une seule clé primaire.
\end{definition}

On distingue deux sortes de clés :
\begin{itemize}
\item une \textbf{clé naturelle} vient du métier : l'adresse e-mail d'un client, par exemple. Mais elle
peut changer, et elle n'est pas toujours unique dans la durée ;
\item une \textbf{clé artificielle} est un numéro sans signification, attribué par la base. C'est le
choix du projet : \code{client_id}, \code{pret_id}, \code{echeance_id}. L'e-mail reste unique grâce à
une contrainte \code{UNIQUE}.
\end{itemize}

La base attribue les numéros grâce aux colonnes \textbf{identité} :
\begin{lstlisting}[style=sql]
client_id   BIGINT GENERATED BY DEFAULT AS IDENTITY PRIMARY KEY  -- numéro automatique, ou fourni
echeance_id BIGINT GENERATED ALWAYS AS IDENTITY PRIMARY KEY      -- numéro toujours automatique
\end{lstlisting}
Avec \code{ALWAYS}, toute tentative d'insérer un numéro à la main est refusée. Avec \code{BY DEFAULT},
on peut fournir le sien : c'est utile au chargement, pour garder le numéro de ligne du fichier source.
L'ancienne syntaxe \code{SERIAL} fait la même chose, en moins strict.

\begin{definition}{Clé étrangère}
Une colonne dont chaque valeur doit exister comme clé primaire dans une autre table. Elle garantit
l'\textbf{intégrité référentielle} : aucune échéance ne peut appartenir à un prêt qui n'existe pas.
\end{definition}

Que faire si l'on supprime la ligne référencée ? C'est l'option \code{ON DELETE} :

\begin{center}\small
\begin{tabularx}{\linewidth}{@{}P{3.3cm} L P{4.3cm}@{}}
\toprule
\textbf{Option} & \textbf{Effet si l'on supprime un client qui a des prêts} & \textbf{Usage} \\
\midrule
\code{RESTRICT} / \code{NO ACTION} & la suppression est \textbf{refusée} & données à conserver : le choix bancaire \\
\code{CASCADE} & les prêts sont supprimés aussi, puis leurs échéances\dots & données sans valeur propre \\
\code{SET NULL} & le \code{client_id} des prêts devient NULL & lien facultatif \\
\bottomrule
\end{tabularx}
\end{center}

\begin{piege}{Piège --- La clé étrangère n'est pas indexée}
PostgreSQL crée un index pour chaque clé primaire et chaque contrainte \code{UNIQUE}, mais
\textbf{pas pour les clés étrangères}. Sans index sur \code{demandes_pret.client_id}, chaque jointure
d'un client vers ses prêts relit toute la table des prêts, et chaque suppression d'un client doit
parcourir toute la table pour vérifier qu'aucun prêt ne le référence. D'où \code{idx_prets_client}.
\end{piege}

Enfin, une contrainte \code{UNIQUE} peut porter sur \textbf{plusieurs colonnes} :
\code{UNIQUE (pret_id, num_echeance)} interdit deux échéances numéro 5 pour le même prêt, tout en
permettant l'échéance numéro 5 de chaque prêt.

\subsection{NOT NULL et CHECK : les règles métier dans la base}

\code{NOT NULL} interdit l'absence de valeur ; \code{DEFAULT} fournit une valeur quand on n'en donne
pas. La contrainte \code{CHECK} accepte une ligne seulement si une condition est vraie :

\begin{lstlisting}[style=sql]
age        SMALLINT     NOT NULL CHECK (age BETWEEN 18 AND 100),
note_risque CHAR(1)     NOT NULL CHECK (note_risque BETWEEN 'A' AND 'G'),
motif_pret VARCHAR(25)  NOT NULL CHECK (motif_pret IN ('etudes', 'sante', 'creation_entreprise',
                                                       'personnel', 'travaux', 'regroupement_dettes')),
\end{lstlisting}

Une contrainte peut porter sur \textbf{plusieurs colonnes de la même ligne}. On lui donne alors un nom,
qui apparaîtra dans les messages d'erreur :
\begin{lstlisting}[style=sql]
CONSTRAINT paiement_borne         CHECK (montant_paye <= montant_du),
CONSTRAINT decision_apres_demande CHECK (date_decision >= date_demande)
\end{lstlisting}

\subsubsection{Deux astuces logiques à connaître}

\textbf{L'équivalence : comparer deux conditions.} En SQL, une condition est une valeur booléenne. On
peut donc écrire \code{(A) = (B)} pour dire « A si et seulement si B ». Exemple : un prêt a un statut
de remboursement \emph{si et seulement si} il est accordé.
\begin{lstlisting}[style=sql]
CONSTRAINT statut_si_accorde CHECK ((statut_dossier = 'accorde') = (statut_pret IS NOT NULL))
\end{lstlisting}

\begin{center}\small
\begin{tabular}{@{}llccc@{}}
\toprule
\code{statut_dossier} & \code{statut_pret} & accordé ? & statut renseigné ? & Ligne acceptée ? \\
\midrule
accorde & sain & vrai & vrai & \ok{} oui \\
refuse & NULL & faux & faux & \ok{} oui \\
accorde & NULL & vrai & faux & \textcolor{rouge}{non} \\
refuse & sain & faux & vrai & \textcolor{rouge}{non} \\
\bottomrule
\end{tabular}
\end{center}

La même astuce relie une décision à la PD et au seuil :
\code{(decision = 'refuser') = (probabilite_defaut >= seuil_decision)}.

\textbf{L'implication : « si A alors B ».} Elle s'écrit \code{NOT A OR B}. Par exemple, « si l'ancienneté
est connue, elle ne dépasse pas l'âge moins 14 ans » :
\begin{lstlisting}[style=sql]
CONSTRAINT anciennete_coherente CHECK (anciennete_emploi IS NULL OR anciennete_emploi <= age - 14)
\end{lstlisting}

\begin{piege}{Piège --- Un CHECK laisse passer NULL}
Une contrainte \code{CHECK} refuse une ligne seulement si la condition est \textbf{fausse}. Si elle
vaut \textbf{NULL} (inconnue), la ligne est acceptée. Or toute comparaison avec NULL donne NULL :
\code{SELECT NULL::int > 0} renvoie NULL, pas faux. Donc \code{CHECK (montant > 0)} accepte un montant
NULL ! Pour l'interdire, il faut ajouter \code{NOT NULL}. C'est voulu pour \code{date_decision} (un
dossier en attente n'a pas encore de décision), pas pour un montant.
\end{piege}

\begin{observation}[Ce qu'un CHECK ne peut pas faire]
Un \code{CHECK} ne voit que la ligne en cours. Il ne peut pas vérifier une règle qui porte sur
\textbf{plusieurs lignes} : « la somme des parts de capital d'un prêt égale le montant prêté », « un
prêt a exactement autant d'échéances que de mois ». Pour ces règles, on écrit des requêtes de contrôle
(exercice 1.4) ou des triggers (chapitre 3).
\end{observation}

\subsection{Les transactions : tout ou rien}

\begin{imagine}
Tu fais un virement de 500~€ de ton compte vers celui d'un ami. La banque doit faire deux opérations :
retirer 500~€ de ton compte, ajouter 500~€ au sien. Si le serveur tombe en panne entre les deux,
l'argent a disparu. Inacceptable : soit les deux opérations ont lieu, soit aucune.
\end{imagine}

\begin{definition}{Transaction et propriétés ACID}
Une transaction est un groupe d'instructions qui forment un tout. PostgreSQL garantit quatre
propriétés :
\begin{itemize}
\item \textbf{Atomicité} : tout est appliqué, ou rien ;
\item \textbf{Cohérence} : la base passe d'un état valide à un état valide (toutes les contraintes sont
respectées) ;
\item \textbf{Isolation} : les transactions simultanées ne voient pas le travail inachevé des autres ;
\item \textbf{Durabilité} : une fois validée, une transaction survit à une panne.
\end{itemize}
\end{definition}

\begin{lstlisting}[style=sql]
BEGIN;                -- début de la transaction
INSERT INTO ...;      -- plusieurs instructions
UPDATE ...;
COMMIT;               -- tout est validé d'un coup (ou ROLLBACK : tout est annulé)
\end{lstlisting}

Sans \code{BEGIN}, chaque instruction est sa propre transaction (mode \emph{autocommit}, celui de psql
et de DBeaver par défaut). Deux usages reviendront tout au long du cours :
\begin{itemize}
\item \textbf{créer des données liées} (un client, son prêt, ses échéances) : si une étape échoue, rien
n'est enregistré ;
\item \textbf{tester sans risque} : \code{BEGIN; ... ; ROLLBACK;} exécute un scénario puis annule tout.
\end{itemize}

\begin{piege}{Piège --- La transaction « interrompue »}
Si une instruction échoue au milieu d'une transaction, PostgreSQL refuse toutes les suivantes avec
le message \emph{current transaction is aborted, commands ignored until end of transaction block}. Il
faut faire \code{ROLLBACK} pour repartir. Dans DBeaver, vérifie aussi le mode de validation (bouton
\emph{Auto-Commit} de la barre d'outils) : en mode manuel, une transaction oubliée garde des verrous et
bloque les autres utilisateurs.
\end{piege}

\subsection{Insérer des données liées : les CTE qui modifient}

Pour créer un prêt, il faut connaître le \code{client_id} que la base vient d'attribuer. La clause
\code{RETURNING} renvoie les lignes insérées, et une \textbf{CTE} (\code{WITH}, détaillée au chapitre 2)
permet de les réutiliser dans l'instruction suivante :

\begin{lstlisting}[style=sql]
WITH nouveau_client AS (
    INSERT INTO clients (...) VALUES (...)
    RETURNING client_id                      -- le numéro attribué par la base
)
INSERT INTO demandes_pret (client_id, ...)
SELECT client_id, ... FROM nouveau_client;   -- on le réutilise sans le connaître à l'avance
\end{lstlisting}

Deux règles à connaître : une CTE qui modifie des données \textbf{s'exécute toujours}, même si personne
ne lit son résultat ; et toutes les parties de l'instruction voient la base \textbf{telle qu'elle était
au début}. Une CTE ne voit pas les lignes insérées par une autre CTE, sauf à travers son
\code{RETURNING}.

\subsection{Le tableau d'amortissement : la dérivation complète}

\begin{situation}{Hands-on --- La mensualité d'Awa}
Awa emprunte $C = 12\,000$~€ sur $n = 36$ mois au taux annuel de 9,5~\%. Elle veut payer \textbf{la même
somme chaque mois}. Quelle est cette mensualité $M$ ?
\end{situation}

\textbf{L'intuition.} Chaque mois, la banque facture des intérêts sur ce qui reste dû : le
\textbf{capital restant dû} (CRD) multiplié par le taux mensuel $r$. La mensualité $M$ paie d'abord ces
intérêts ; le reste rembourse le capital. Au début, le CRD est élevé, donc les intérêts aussi : la part
de capital est petite. Puis le CRD baisse, les intérêts aussi, et la part de capital grandit
(figure~\ref{fig:amortissement}).

\textbf{La dérivation.} Le taux mensuel est $r = \dfrac{\text{taux annuel}}{12 \times 100}$. Notons
$CRD_k$ le capital restant dû après $k$ mensualités.
\begin{enumerate}
\item Au départ, $CRD_0 = C$.
\item Chaque mois, le CRD produit des intérêts puis la mensualité est payée :
\[ CRD_k = CRD_{k-1}\,(1+r) - M \]
\item En déroulant cette relation $k$ fois, on obtient une somme géométrique :
\[ CRD_k = C(1+r)^k - M\left[(1+r)^{k-1} + \dots + (1+r) + 1\right]
         = C(1+r)^k - M\,\frac{(1+r)^k - 1}{r} \]
\item Le prêt est remboursé après $n$ mois : $CRD_n = 0$. On isole $M$ :
\[ M\,\frac{(1+r)^n - 1}{r} = C(1+r)^n \quad\Longrightarrow\quad
   \boxed{\;M = \frac{C \, r}{1 - (1+r)^{-n}}\;} \]
\end{enumerate}

\begin{calcul}{la mensualité d'Awa}
\begin{enumerate}
\item Taux mensuel : $r = 9{,}5 / 1200 = 0{,}0079167$.
\item $(1+r)^{36} = 1{,}328271$, donc $(1+r)^{-36} = 0{,}752859$ et $1 - (1+r)^{-36} = 0{,}247141$.
\item $C \times r = 12\,000 \times 0{,}0079167 = 95{,}00$.
\item $M = 95{,}00 / 0{,}247141 = 384{,}395$, arrondi au centime : $\mathbf{384{,}40~€}$.
\end{enumerate}
Premier mois : intérêts $= 12\,000 \times r = 95{,}00$~€, capital remboursé $= 384{,}40 - 95{,}00 =
289{,}40$~€, nouveau CRD $= 11\,710{,}60$~€. Coût total du crédit : 1\,838,21~€.
\end{calcul}

\begin{figure}[ht]
\centering
\begin{tikzpicture}
\begin{axis}[ybar stacked, bar width=7pt, width=15cm, height=6.2cm, ymin=0, ymax=420,
  xmin=0, xmax=37, xlabel={Numéro de la mensualité}, ylabel={Euros},
  legend style={at={(0.5,1.03)}, anchor=south, legend columns=2, draw=none, font=\small},
  xtick={1,6,12,18,24,30,36}, tick label style={font=\footnotesize}, label style={font=\small},
  grid=major, grid style={gray!20}, axis lines*=left]
\addplot[fill=bleuML!75, draw=none] coordinates {
(1,289.40) (2,291.69) (3,294.00) (4,296.33) (5,298.67) (6,301.04) (7,303.42) (8,305.83) (9,308.24)
(10,310.69) (11,313.14) (12,315.62) (13,318.13) (14,320.64) (15,323.18) (16,325.73) (17,328.32)
(18,330.92) (19,333.53) (20,336.18) (21,338.83) (22,341.52) (23,344.23) (24,346.95) (25,349.69)
(26,352.46) (27,355.26) (28,358.06) (29,360.90) (30,363.76) (31,366.64) (32,369.54) (33,372.47)
(34,375.41) (35,378.39) (36,381.19)};
\addplot[fill=orangeEx!80, draw=none] coordinates {
(1,95.00) (2,92.71) (3,90.40) (4,88.07) (5,85.73) (6,83.36) (7,80.98) (8,78.57) (9,76.16) (10,73.71)
(11,71.26) (12,68.78) (13,66.27) (14,63.76) (15,61.22) (16,58.67) (17,56.08) (18,53.48) (19,50.87)
(20,48.22) (21,45.57) (22,42.88) (23,40.17) (24,37.45) (25,34.71) (26,31.94) (27,29.14) (28,26.34)
(29,23.50) (30,20.64) (31,17.76) (32,14.86) (33,11.93) (34,8.99) (35,6.01) (36,3.02)};
\legend{Part de capital, Part d'intérêts}
\end{axis}
\end{tikzpicture}
\caption{Le tableau d'amortissement d'Awa, tel que ta base le génère : la mensualité reste à
384,40~€, mais sa composition change. La dernière (384,21~€) solde exactement le capital restant.}
\label{fig:amortissement}
\end{figure}

\begin{piege}{Piège --- Les arrondis au centime}
La mensualité est arrondie au centime (384,40 au lieu de 384,395\dots). Si l'on arrondit ensuite
\emph{chaque} colonne indépendamment (intérêts, capital, CRD), le CRD d'un mois ne vaut plus exactement
le CRD du mois précédent moins la part de capital : des écarts d'un centime apparaissent. La bonne
méthode : arrondir les CRD, puis définir la part de capital comme la \textbf{différence} de deux CRD
arrondis, et ajuster la dernière mensualité pour solder exactement le prêt. Les contrôles de
l'exercice 1.4 vérifient tout cela.
\end{piege}

\subsection{Générer des lignes : generate\_series}

Comment produire 36 lignes à partir d'un seul prêt ? La fonction \code{generate_series(a, b)} renvoie
une table d'entiers de $a$ à $b$ :
\begin{lstlisting}[style=sql]
SELECT p.pret_id, k AS num_echeance,
       (p.date_premiere_echeance + (k - 1) * INTERVAL '1 month')::date AS date_echeance
FROM demandes_pret p
CROSS JOIN generate_series(1, p.duree_mois) AS k     -- k = 1, 2, ..., duree_mois
WHERE p.pret_id = 32582;
\end{lstlisting}
Le \code{CROSS JOIN} associe chaque prêt à chaque valeur de $k$ : un prêt de 36 mois produit 36 lignes.
La fonction peut utiliser les colonnes du prêt (\code{p.duree_mois}) : PostgreSQL l'évalue pour chaque
prêt.

\subsection{Les index B-Tree : retrouver une ligne parmi un million}

\begin{analogie}{l'annuaire téléphonique}
Pour trouver le numéro de « Koné » dans un annuaire, tu ne lis pas toutes les pages : tu l'ouvres au
milieu, tu vois que tu es à « M », tu reviens en arrière, et en quelques gestes tu es à la bonne page.
C'est possible parce que l'annuaire est \textbf{trié}. Un index est un annuaire de la table : une copie
triée d'une ou plusieurs colonnes, avec pour chaque valeur l'adresse de la ligne.
\end{analogie}

L'index B-Tree (\emph{balanced tree}, l'index par défaut de PostgreSQL) range les valeurs dans un arbre
équilibré. Chaque nœud contient des centaines de valeurs : trois ou quatre niveaux suffisent pour
plus d'un million de lignes (figure~\ref{fig:btree}).

\begin{figure}[ht]
\centering
\begin{tikzpicture}[font=\footnotesize, >=Stealth,
  noeud/.style={draw=grisTexte, fill=white, minimum height=0.6cm, inner sep=3pt},
  chemin/.style={draw=bleuML, fill=bleuFond, line width=0.9pt}]
  \node[noeud, chemin] (r) at (0,0) {racine : \;\dots\; | 8\,200 | 16\,400 | 24\,600 | \dots};
  \node[noeud, chemin] (i1) at (-4.2,-1.5) {\dots\; | 40 | 82 | 125 | \dots};
  \node[noeud] (i2) at (0,-1.5) {\dots\; | 8\,240 | 8\,281 | \dots};
  \node[noeud] (i3) at (4.2,-1.5) {\dots\; | 16\,440 | 16\,481 | \dots};
  \node[noeud] (f0) at (-6.3,-3.1) {38 | 39 | 40};
  \node[noeud, chemin] (f1) at (-3.6,-3.1) {\textbf{61} (18 août) | 61 (15 oct.) | \dots};
  \node[noeud] (f2) at (-0.5,-3.1) {82 | 83 | \dots};
  \node[draw=grisTitre, dashed, minimum width=3.3cm, minimum height=0.7cm, align=center] (t) at (2.7,-3.1)
     {table : ligne de l'échéance};
  \draw[->, chemin] (r) -- (i1);
  \draw[->, grisTexte] (r) -- (i2);
  \draw[->, grisTexte] (r) -- (i3);
  \draw[->, grisTexte] (i1) -- (f0);
  \draw[->, chemin] (i1) -- (f1);
  \draw[->, grisTexte] (i1) -- (f2);
  \draw[->, chemin, dashed] (f1.south) to[bend right=18] node[below, font=\scriptsize] {adresse de la ligne} (t.south);
\end{tikzpicture}
\caption{Recherche de l'impayé le plus ancien du prêt 61 dans l'index partiel
\code{idx_echeances_impayees (pret_id, date_echeance)} : trois nœuds lus au lieu de 1,1 million de lignes.}
\label{fig:btree}
\end{figure}

\subsubsection{Ce que ça change, mesuré sur ta base}

La commande \code{EXPLAIN ANALYZE} exécute une requête et affiche le \textbf{plan} choisi par
PostgreSQL avec les temps réels. Voici le calcul du retard du prêt 61 :

\begin{lstlisting}[style=sql]
EXPLAIN ANALYZE
SELECT coalesce(DATE '2026-09-30' - min(date_echeance), 0)
FROM echeancier_remboursements
WHERE pret_id = 61 AND date_echeance <= DATE '2026-09-30' AND montant_paye < montant_du;
\end{lstlisting}

\textbf{Avec l'index partiel} : PostgreSQL descend l'arbre et lit une seule entrée.
\begin{lstlisting}[style=sortie]
->  Index Only Scan using idx_echeances_impayees on echeancier_remboursements
      Index Cond: ((pret_id = 61) AND (date_echeance <= '2026-09-30'::date))
Execution Time: 0.318 ms
\end{lstlisting}

\textbf{Sans index} (supprimés le temps du test, dans une transaction annulée) : PostgreSQL lit toute
la table, en parallèle sur trois processus.
\begin{lstlisting}[style=sortie]
->  Parallel Seq Scan on echeancier_remboursements
      Filter: ((date_echeance <= '2026-09-30'::date) AND (montant_paye < montant_du) AND (pret_id = 61))
      Rows Removed by Filter: 369324
Execution Time: 113.944 ms
\end{lstlisting}

Soit un facteur \textbf{350}. Pour un prêt, 114~ms paraissent acceptables. Mais au chapitre 3, une
fonction calculera le retard des 32\,410 prêts un par un : avec l'index, 0,25 seconde ; sans index,
plus de 20 minutes.

\subsubsection{Les règles à connaître}

\begin{itemize}
\item \textbf{Index composite et préfixe gauche.} L'index \code{(pret_id, date_echeance)} est trié par
prêt, puis par date à l'intérieur de chaque prêt, comme un annuaire trié par nom puis prénom. Il sert
pour \code{WHERE pret_id = 61} et pour \code{WHERE pret_id = 61 AND date_echeance <= ...}, mais
\textbf{pas} pour \code{WHERE date_echeance <= ...} seul.
\item \textbf{Index partiel.} \code{WHERE montant_paye < montant_du} n'indexe que les échéances non
soldées : 13~Mo au lieu de 33~Mo pour un index complet, et il reste à jour automatiquement (une échéance
payée en sort d'elle-même).
\item \textbf{Un index a un coût.} Chaque \code{INSERT} ou \code{UPDATE} doit aussi mettre à jour les
index, et ils occupent de la place. On n'indexe que les colonnes réellement utilisées pour filtrer ou
joindre.
\item \textbf{\code{ANALYZE}} met à jour les statistiques que PostgreSQL utilise pour choisir son plan :
à lancer après un gros chargement (le chargeur le fait).
\end{itemize}

\begin{piege}{Piège --- Une fonction sur la colonne indexée}
\code{WHERE date_trunc('month', date_echeance) = DATE '2026-09-01'} ne peut pas utiliser un index sur
\code{date_echeance} : l'index est trié sur la date, pas sur le résultat de la fonction. Écris plutôt
un intervalle : \code{WHERE date_echeance >= DATE '2026-09-01' AND date_echeance < DATE '2026-10-01'}.
\end{piege}

% =============================================================================
\section{Exercices guidés}

\begin{exercice}{1.1}{Ouvrir un nouveau dossier de prêt complet, en une transaction}
Crée le dossier d'Awa Koné, en \textbf{une seule transaction} :
\begin{enumerate}
\item la cliente : Awa Koné, \code{awa.kone@exemple.fr}, 29 ans, 42\,000~€ de revenu, locataire,
4 ans d'ancienneté dans l'emploi, 6 ans d'historique de crédit, aucun défaut antérieur, inscrite le
15/09/2026 ;
\item son prêt : 12\,000~€ sur 36 mois à 9,5~\%, motif études, note B, demandé le 20/09/2026, accordé
le 25/09/2026, première échéance le 25/10/2026, statut « sain » ;
\item la prédiction du modèle champion \code{xgboost_monotone_v2} qui a justifié la décision : PD
0,118, seuil 0,255, décision « accorder », variable SHAP principale \code{taux_effort} (0,912), calculée
le 24/09/2026 à 14~h~30 ;
\item ses 36 échéances, à mensualité constante.
\end{enumerate}
Aucun identifiant ne doit être écrit en dur : utilise ceux que la base attribue. Vérifie ensuite le
nombre d'échéances, le capital total remboursé et le coût du crédit.
\end{exercice}

\begin{correction}{1.1}
\textbf{Étape 1 --- Ouvrir la transaction.} \code{BEGIN} : si une seule des quatre insertions échoue,
aucune ne sera enregistrée.

\textbf{Étape 2 --- La cliente.} Une CTE insère la cliente et renvoie son \code{client_id}.

\textbf{Étape 3 --- Le prêt.} La CTE suivante lit ce \code{client_id} dans \code{nouveau_client}. Son
\code{RETURNING} renvoie tout ce dont l'échéancier aura besoin, sous des noms courts : le capital $c$,
la durée $n$, le taux mensuel $r$ et la date de première échéance.

\textbf{Étape 4 --- La prédiction.} Une CTE insère le score. Personne ne la lit ensuite : elle
s'exécute quand même.

\textbf{Étape 5 --- Les mensualités.} \code{generate_series(1, n)} produit les 36 numéros ; la mensualité
est calculée par la formule d'annuité et arrondie au centime ; chaque date part de la première
échéance (voir le piège des mois).

\textbf{Étape 6 --- Le capital restant dû.} Pour chaque $k$, on calcule $CRD_{k-1}$ et $CRD_k$ par la
formule fermée, arrondis au centime.

\textbf{Étape 7 --- L'insertion.} Pour les 35 premières échéances : montant $= M$, capital
$= CRD_{k-1} - CRD_k$, intérêts $= M -$ capital. La dernière solde exactement le capital restant :
capital $= CRD_{35}$, intérêts $= CRD_{35} \times r$.

\begin{lstlisting}[style=long]
BEGIN;

WITH nouveau_client AS (
    INSERT INTO clients (nom_complet, email, age, revenu_annuel, statut_logement,
                         anciennete_emploi, anciennete_credit, defaut_anterieur, date_inscription)
    VALUES ('Awa Koné', 'awa.kone@exemple.fr', 29, 42000, 'locataire', 4, 6, false, DATE '2026-09-15')
    RETURNING client_id
),
nouveau_pret AS (
    INSERT INTO demandes_pret (client_id, montant_demande, duree_mois, taux_interet, motif_pret, note_risque,
                               date_demande, date_decision, date_premiere_echeance, statut_dossier, statut_pret)
    SELECT client_id, 12000, 36, 9.5, 'etudes', 'B',
           DATE '2026-09-20', DATE '2026-09-25', DATE '2026-10-25', 'accorde', 'sain'
    FROM nouveau_client
    RETURNING pret_id, montant_demande AS c, duree_mois AS n, taux_interet / 1200 AS r, date_premiere_echeance
),
-- Le score du modèle champion qui a justifié la décision (PD 11,8 % < seuil 25,5 % : accordé).
-- Une CTE qui modifie des données s'exécute même si personne ne la lit ensuite.
score AS (
    INSERT INTO predictions_scoring_ml (pret_id, version_modele, est_champion, probabilite_defaut,
                                        seuil_decision, decision, shap_top_feature, shap_top_valeur, date_prediction)
    SELECT pret_id, 'xgboost_monotone_v2', true, 0.118, 0.255, 'accorder', 'taux_effort', 0.912,
           TIMESTAMP '2026-09-24 14:30'
    FROM nouveau_pret
),
-- Mensualité constante : M = C × r / (1 − (1 + r)^−n)
mensualites AS (
    SELECT p.pret_id, p.c, p.r, p.n, k AS num_echeance,
           (p.date_premiere_echeance + (k - 1) * INTERVAL '1 month')::DATE AS date_echeance,
           round(p.c * p.r / (1 - (1 + p.r) ^ (-p.n)), 2)                    AS mensualite
    FROM nouveau_pret p
    CROSS JOIN generate_series(1, p.n) AS k                   -- une ligne par mensualité
),
-- Capital restant dû après k mensualités : CRD(k) = C(1 + r)^k − M((1 + r)^k − 1) / r, arrondi au centime.
-- La part de capital est la DIFFÉRENCE de deux CRD arrondis : le CRD baisse exactement de la part de capital.
calcul AS (
    SELECT pret_id, num_echeance, date_echeance, mensualite, r, n,
           round(c * (1 + r) ^ (num_echeance - 1) - mensualite * ((1 + r) ^ (num_echeance - 1) - 1) / r, 2) AS crd_avant,
           round(c * (1 + r) ^ num_echeance       - mensualite * ((1 + r) ^ num_echeance - 1) / r, 2)       AS crd_apres
    FROM mensualites
)
INSERT INTO echeancier_remboursements (pret_id, num_echeance, date_echeance, montant_du,
                                       part_capital, part_interets, capital_restant_du)
SELECT pret_id, num_echeance, date_echeance,
       -- la dernière échéance solde exactement le capital restant
       CASE WHEN num_echeance = n THEN crd_avant + round(crd_avant * r, 2) ELSE mensualite END,
       CASE WHEN num_echeance = n THEN crd_avant ELSE crd_avant - crd_apres END,
       CASE WHEN num_echeance = n THEN round(crd_avant * r, 2) ELSE mensualite - (crd_avant - crd_apres) END,
       CASE WHEN num_echeance = n THEN 0 ELSE crd_apres END
FROM calcul;

COMMIT;
\end{lstlisting}

\textbf{Étape 8 --- Vérifier.} On retrouve le prêt par l'e-mail de la cliente, sans connaître son
numéro. \code{FILTER} (vu au chapitre 2) isole une échéance dans un agrégat.

\begin{lstlisting}[style=long]
SELECT p.pret_id, count(*) AS nb_echeances, sum(e.part_capital) AS capital_rembourse,
       sum(e.part_interets) AS cout_du_credit,
       max(e.montant_du) FILTER (WHERE e.num_echeance = 1) AS mensualite,
       max(e.montant_du) FILTER (WHERE e.num_echeance = p.duree_mois) AS derniere_echeance
FROM demandes_pret p JOIN echeancier_remboursements e USING (pret_id)
WHERE p.client_id = (SELECT client_id FROM clients WHERE email = 'awa.kone@exemple.fr')
GROUP BY p.pret_id;
\end{lstlisting}

\begin{center}\small
\begin{tabular}{@{}rrrrrr@{}}
\toprule
\code{pret_id} & \code{nb_echeances} & \code{capital_rembourse} & \code{cout_du_credit} & \code{mensualite} & \code{derniere_echeance} \\
\midrule
32582 & 36 & 12\,000,00 & 1\,838,21 & 384,40 & 384,21 \\
\bottomrule
\end{tabular}
\end{center}

\textbf{À retenir.} Les identifiants circulent par \code{RETURNING} au lieu d'être écrits en dur ; la
transaction garantit le « tout ou rien » ; l'échéancier est calculé en une seule requête ensembliste,
sans boucle.
\end{correction}

\begin{exercice}{1.2}{Les contraintes refusent les données incohérentes}
Écris quatre instructions qui violent chacune une règle, exécute-les \textbf{une par une} et lis le
message d'erreur :
\begin{enumerate}[label=\alph*)]
\item créer un client de 15 ans ;
\item enregistrer sur une échéance future impayée un paiement supérieur au montant dû ;
\item supprimer le client 2, qui a un prêt ;
\item créer un prêt refusé dont le statut de remboursement est « sain ».
\end{enumerate}
\end{exercice}

\begin{correction}{1.2}
\begin{lstlisting}[style=long]
-- a) Un client de 15 ans
INSERT INTO clients (nom_complet, email, age, revenu_annuel, statut_logement, anciennete_credit,
                     defaut_anterieur, date_inscription)
VALUES ('Test', 'test@exemple.fr', 15, 20000, 'locataire', 0, false, CURRENT_DATE);

-- b) Payer plus que le montant dû, sur une échéance future encore impayée
UPDATE echeancier_remboursements SET montant_paye = montant_du + 1
WHERE echeance_id = (SELECT min(echeance_id) FROM echeancier_remboursements
                     WHERE montant_paye = 0 AND date_echeance > DATE '2026-09-30');

-- c) Supprimer un client qui a des prêts
DELETE FROM clients WHERE client_id = 2;

-- d) Un prêt refusé avec un statut de remboursement
INSERT INTO demandes_pret (client_id, montant_demande, duree_mois, taux_interet, motif_pret, note_risque,
                           date_demande, statut_dossier, statut_pret)
VALUES (2, 1000, 12, 10, 'etudes', 'A', DATE '2026-09-01', 'refuse', 'sain');
\end{lstlisting}

Les quatre messages obtenus :
\begin{lstlisting}[style=sortie]
a) ERROR: new row for relation "clients" violates check constraint "clients_age_check"
b) ERROR: new row for relation "echeancier_remboursements" violates check constraint "paiement_borne"
c) ERROR: update or delete on table "clients" violates foreign key constraint
          "demandes_pret_client_id_fkey" on table "demandes_pret"
   DETAIL: Key (client_id)=(2) is still referenced from table "demandes_pret".
d) ERROR: new row for relation "demandes_pret" violates check constraint "statut_si_accorde"
\end{lstlisting}

\textbf{Lecture des messages.} En a), la contrainte n'a pas de nom explicite : PostgreSQL l'a nommée
\code{clients_age_check}. En b), d) le nom choisi dans le schéma rend l'erreur immédiatement
compréhensible : c'est l'intérêt de nommer ses contraintes.

\textbf{Pourquoi une échéance future en b) ?} Sur une échéance déjà payée, le nouveau montant violerait
\emph{deux} contraintes : \code{paiement_borne}, mais aussi \code{date_si_solde} (une date de paiement
alors que le montant payé ne vaut plus le montant dû). PostgreSQL signale alors la première qu'il
vérifie, \code{date_si_solde}, ce qui brouille la démonstration.

\textbf{Pourquoi la base, et pas l'application ?} Une application peut avoir un bug, un analyste peut
écrire une requête directement, un script de reprise peut contourner l'interface. La contrainte en
base protège \textbf{tous} les chemins d'accès.
\end{correction}

% =============================================================================
\section{Exercices à faire seul}

\begin{exercice}{1.3}{Les deux tables du chapitre 3}
Le chapitre 3 a besoin de deux nouvelles tables. Écris leur \code{CREATE TABLE}, avec les types et
toutes les contraintes qui te paraissent nécessaires.

\textbf{1. \code{journal_statuts}} --- le journal d'audit des changements de statut d'un prêt. Chaque
ligne enregistre : un identifiant automatique ; le prêt concerné ; l'ancien statut (qui peut être
absent) et le nouveau (obligatoire) ; le nombre de jours de retard au moment du changement (jamais
négatif) ; la date d'effet ; l'horodatage de l'enregistrement (par défaut, l'instant présent) ;
l'utilisateur qui a fait le changement (par défaut, l'utilisateur connecté). Une ligne n'a de sens que
si le statut a réellement changé. Pense à l'index utile.

\textbf{2. \code{provisions_mensuelles}} --- une « photographie » du risque de chaque prêt à chaque
date d'arrêté : le retard, le stage IFRS~9 (1, 2 ou 3), l'EAD, la PD utilisée, la LGD et l'ECL. Un prêt
ne peut avoir qu'\textbf{une ligne par arrêté}.
\end{exercice}

\begin{indice}
\begin{itemize}
\item « Ancien différent de nouveau » : que donne \code{NULL <> 'sain'} ? Cherche l'opérateur
\code{IS DISTINCT FROM}.
\item L'utilisateur connecté s'obtient avec \code{session_user} ; l'instant présent avec \code{now()}.
\item « Une ligne par prêt et par arrêté » : une clé primaire peut porter sur deux colonnes.
\item Les clés étrangères : quelle option \code{ON DELETE} pour un journal d'audit ?
\end{itemize}
\end{indice}

\begin{exercice}{1.4}{Contrôles qualité de l'échéancier}
Les contraintes vérifient chaque \emph{ligne}. Certaines règles portent sur un prêt \emph{entier} ou sur
deux échéances successives : il faut les contrôler par des requêtes. Écris des requêtes qui comptent les
anomalies suivantes (le résultat attendu est 0 partout) :
\begin{enumerate}[label=\alph*)]
\item par prêt accordé : la somme des parts de capital diffère du montant prêté ; le nombre
d'échéances ou le dernier numéro diffère de la durée ; le CRD de la dernière échéance n'est pas nul ;
une échéance a un montant dû différent de capital + intérêts ;
\item entre deux échéances successives d'un même prêt : le CRD ne baisse pas exactement de la part
de capital ; l'écart entre les deux dates n'est pas compris entre 28 et 31 jours ;
\item un paiement est daté d'avant l'octroi du prêt.
\end{enumerate}
\end{exercice}

\begin{indice}
\begin{itemize}
\item a) Un \code{GROUP BY p.pret_id} sur la jointure prêts--échéances, puis compter les prêts en
anomalie. \code{count(*) FILTER (WHERE condition)} compte les lignes qui vérifient la condition (vu au
chapitre 2) ; \code{sum(CASE WHEN condition THEN 1 ELSE 0 END)} fait la même chose.
\item b) Joins la table à elle-même (\textbf{auto-jointure}) : l'échéance \code{prec} et l'échéance
\code{suiv} du même prêt, avec \code{suiv.num_echeance = prec.num_echeance + 1}.
\item Si tu trouves des écarts d'un centime, relis le piège des arrondis.
\end{itemize}
\end{indice}

% =============================================================================
\section{Synthèse pour les entretiens}

\begin{resume}
\begin{enumerate}
\item \textbf{Normaliser} : chaque fait stocké une seule fois. Les tables d'analyse pour le ML peuvent
être dénormalisées, pas le système source.
\item \textbf{Types} : \code{NUMERIC} pour l'argent ; \code{DATE - DATE} donne des jours ; toujours
ajouter des mois à partir de la première échéance.
\item \textbf{Clés} : primaire (unique et non nulle), étrangère (intégrité référentielle, \code{RESTRICT}
en banque), \code{UNIQUE} composite. Les clés étrangères ne sont pas indexées automatiquement.
\item \textbf{CHECK} : règles sur une ligne ; équivalence \code{(A) = (B)}, implication \code{NOT A OR B} ;
un CHECK laisse passer NULL.
\item \textbf{Transactions} : ACID, tout ou rien ; \code{BEGIN ... ROLLBACK} pour tester sans risque.
\item \textbf{CTE et RETURNING} pour insérer des données liées sans écrire d'identifiant en dur.
\item \textbf{Mensualité} : $M = C\,r / (1 - (1+r)^{-n})$ ; attention aux arrondis au centime.
\item \textbf{Index B-Tree} : recherche en quelques nœuds ; préfixe gauche ; index partiel ; mesuré :
0,3~ms contre 114~ms.
\end{enumerate}
\end{resume}

\begin{entretien}
\begin{itemize}
\item « J'ai conçu un schéma normalisé de six tables, plus deux tables d'audit et de provisions. Les
règles métier sont dans la base : des contraintes \code{CHECK} nommées, y compris entre colonnes, comme
l'équivalence entre une décision de refus et une PD au-dessus du seuil. »
\item « J'utilise \code{NUMERIC} pour les montants, parce qu'un \code{FLOAT} ne représente pas
exactement 0,10 et que les totaux comptables doivent tomber juste au centime. »
\item « Toutes les clés étrangères sont en \code{ON DELETE RESTRICT} : supprimer un client ne doit
jamais effacer l'historique de ses prêts. Et je les indexe, car PostgreSQL ne le fait pas. »
\item « Le calcul du retard utilise un index partiel sur les seules échéances impayées : 13~Mo au lieu
de 33, et une requête en 0,3~ms au lieu de 114. »
\end{itemize}
\end{entretien}

\subsection*{Questions fréquentes en entretien}

\begin{description}[leftmargin=0pt, itemsep=0.5em]
\item[Clé primaire ou contrainte UNIQUE ?] Une table a une seule clé primaire, jamais NULL ; elle peut
avoir plusieurs contraintes \code{UNIQUE}, qui acceptent les NULL (plusieurs NULL ne sont pas
considérés comme égaux).
\item[DELETE, TRUNCATE ou DROP ?] \code{DELETE} supprime des lignes (avec un \code{WHERE}, ligne à
ligne, triggers déclenchés) ; \code{TRUNCATE} vide toute la table d'un coup, très vite ; \code{DROP}
supprime la table elle-même.
\item[Pourquoi ne pas tout mettre en CASCADE ?] Parce qu'une seule suppression pourrait effacer en
chaîne des milliers de lignes d'historique, sans avertissement.
\item[Quand un index ralentit-il ?] Il ralentit les écritures (chaque insertion met à jour chaque
index) et devient inutile si la requête lit une grande partie de la table.
\item[Qu'est-ce qu'une transaction ?] Un groupe d'instructions atomique : tout ou rien. Avec les
propriétés ACID.
\end{description}
